% ==================== THEME 4 — BOUCLES SIMPLES ====================

% ---- boucle-for-trace ----
\element{bank-boucle-for-trace}{
    \begin{question}{boucle-for-trace-v1}
        Soit le code ci-dessous, qu'affichera le programme en fin d'exécution ?
        \lstinputlisting{sources/codes/boucle-for-1.c}
        \begin{multicols}{6}
            \begin{reponses}
                \bonne{10}
                \mauvaise{15}
                \mauvaise{0}
                \mauvaise{5}
                \mauvaise{4}
                \mauvaise{i}
            \end{reponses}
        \end{multicols}
    \end{question}
}
\element{bank-boucle-for-trace}{
    \begin{question}{boucle-for-trace-v2}
        Soit le code ci-dessous, qu'affichera le programme en fin d'exécution ?
        \lstinputlisting{sources/codes/boucle-for-2.c}
        \begin{multicols}{6}
            \begin{reponses}
                \bonne{6}
                \mauvaise{10}
                \mauvaise{0}
                \mauvaise{4}
                \mauvaise{3}
                \mauvaise{i}
            \end{reponses}
        \end{multicols}
    \end{question}
}
\element{bank-boucle-for-trace}{
    \begin{question}{boucle-for-trace-v3}
        Soit le code ci-dessous, qu'affichera le programme en fin d'exécution ?
        \lstinputlisting{sources/codes/boucle-for-3.c}
        \begin{multicols}{6}
            \begin{reponses}
                \bonne{15}
                \mauvaise{21}
                \mauvaise{0}
                \mauvaise{6}
                \mauvaise{5}
                \mauvaise{i}
            \end{reponses}
        \end{multicols}
    \end{question}
}
\element{boucles}{
    \restituegroupe[1]{bank-boucle-for-trace}
}

% ---- boucle-while-trace ----
\element{bank-boucle-while-trace}{
    \begin{question}{boucle-while-trace-v1}
        Soit le code ci-dessous, qu'affichera le programme en fin d'exécution ?
        \lstinputlisting{sources/codes/boucle-while-1.c}
        \begin{multicols}{5}
            \begin{reponses}
                \bonne{0123}
                \mauvaise{1234}
                \mauvaise{01234}
                \mauvaise{4}
                \mauvaise{Boucle infinie}
            \end{reponses}
        \end{multicols}
    \end{question}
}
\element{bank-boucle-while-trace}{
    \begin{question}{boucle-while-trace-v2}
        Soit le code ci-dessous, qu'affichera le programme en fin d'exécution ?
        \lstinputlisting{sources/codes/boucle-while-2.c}
        \begin{multicols}{5}
            \begin{reponses}
                \bonne{01234}
                \mauvaise{12345}
                \mauvaise{012345}
                \mauvaise{5}
                \mauvaise{Boucle infinie}
            \end{reponses}
        \end{multicols}
    \end{question}
}
\element{bank-boucle-while-trace}{
    \begin{question}{boucle-while-trace-v3}
        Soit le code ci-dessous, qu'affichera le programme en fin d'exécution ?
        \lstinputlisting{sources/codes/boucle-while-3.c}
        \begin{multicols}{5}
            \begin{reponses}
                \bonne{012}
                \mauvaise{123}
                \mauvaise{0123}
                \mauvaise{3}
                \mauvaise{Boucle infinie}
            \end{reponses}
        \end{multicols}
    \end{question}
}
\element{boucles}{
    \restituegroupe[1]{bank-boucle-while-trace}
}

% ---- boucle-for-vs-while ----
\element{bank-boucle-for-vs-while}{
    \begin{question}{boucle-for-vs-while-v1}
        Quelle boucle utiliser lorsque l'on connaît \textbf{à l'avance} le nombre exact de répétitions ?
        \begin{multicols}{3}
            \begin{reponses}
                \bonne{\lstinline[language=C]|for|}
                \mauvaise{\lstinline[language=C]|while| uniquement}
                \mauvaise{\lstinline[language=C]|if|}
            \end{reponses}
        \end{multicols}
    \end{question}
}
\element{bank-boucle-for-vs-while}{
    \begin{question}{boucle-for-vs-while-v2}
        Quelle boucle est la plus adaptée pour répéter un traitement un nombre \textbf{fixe et connu} de fois (par exemple 20 fois) ?
        \begin{multicols}{3}
            \begin{reponses}
                \bonne{\lstinline[language=C]|for|}
                \mauvaise{\lstinline[language=C]|while| uniquement}
                \mauvaise{\lstinline[language=C]|if|}
            \end{reponses}
        \end{multicols}
    \end{question}
}
\element{bank-boucle-for-vs-while}{
    \begin{question}{boucle-for-vs-while-v3}
        Pour parcourir un intervalle dont les bornes sont connues dès le départ, quelle structure de boucle est la plus naturelle en C ?
        \begin{multicols}{3}
            \begin{reponses}
                \bonne{\lstinline[language=C]|for|}
                \mauvaise{\lstinline[language=C]|while| uniquement}
                \mauvaise{\lstinline[language=C]|if|}
            \end{reponses}
        \end{multicols}
    \end{question}
}
\element{boucles}{
    \restituegroupe[1]{bank-boucle-for-vs-while}
}

% ---- boucle-condition-arret ----
\element{bank-boucle-condition-arret}{
    \begin{question}{boucle-condition-arret-v1}
        Quel risque court le programme suivant ?
        \lstinputlisting{sources/codes/boucle-infinie.c}
        \begin{multicols}{2}
            \begin{reponses}
                \bonne{Boucle infinie, car \texttt{i} n'est jamais incrémenté}
                \mauvaise{Erreur de compilation}
                \mauvaise{Le programme n'affiche rien}
                \mauvaise{Le programme s'arrête après un seul tour}
            \end{reponses}
        \end{multicols}
    \end{question}
}
\element{bank-boucle-condition-arret}{
    \begin{question}{boucle-condition-arret-v2}
        Quel risque court le programme suivant ?
        \lstinputlisting{sources/codes/boucle-infinie-2.c}
        \begin{multicols}{2}
            \begin{reponses}
                \bonne{Boucle infinie, car \texttt{n} n'est jamais incrémenté}
                \mauvaise{Erreur de compilation}
                \mauvaise{Le programme n'affiche rien}
                \mauvaise{Le programme s'arrête après un seul tour}
            \end{reponses}
        \end{multicols}
    \end{question}
}
\element{bank-boucle-condition-arret}{
    \begin{question}{boucle-condition-arret-v3}
        Quel risque court le programme suivant ?
        \lstinputlisting{sources/codes/boucle-infinie-3.c}
        \begin{multicols}{2}
            \begin{reponses}
                \bonne{Boucle infinie, car \texttt{k} n'est jamais décrémenté}
                \mauvaise{Erreur de compilation}
                \mauvaise{Le programme n'affiche rien}
                \mauvaise{Le programme s'arrête après un seul tour}
            \end{reponses}
        \end{multicols}
    \end{question}
}
\element{boucles}{
    \restituegroupe[1]{bank-boucle-condition-arret}
}

% ---- boucle-offbyone ----
% Durci : le simple off-by-one sur une borne unique (\texttt{i <= N}) est
% acquis en fin de séquence. On ajoute une seconde condition combinée par
% \texttt{\&\&} qui provoque un arrêt anticipé de la boucle : le piège n'est
% plus de compter les bornes mais de repérer QUELLE condition arrête la
% boucle en premier.
\element{bank-boucle-offbyone}{
    \begin{question}{boucle-offbyone-v1}
        Combien de fois le message \texttt{"Bonjour"} est-il affiché par \lstinline[language=C]|for (int i = 1; i <= 10 && i != 6; i++) { printf("Bonjour"); }| ?
        \begin{multicols}{5}
            \begin{reponses}
                \mauvaise{10}
                \bonne{5}
                \mauvaise{6}
                \mauvaise{4}
                \mauvaise{Boucle infinie}
            \end{reponses}
        \end{multicols}
    \end{question}
}
\element{bank-boucle-offbyone}{
    \begin{question}{boucle-offbyone-v2}
        Combien de fois le message \texttt{"Bonjour"} est-il affiché par \lstinline[language=C]|for (int i = 1; i <= 8 && i != 5; i++) { printf("Bonjour"); }| ?
        \begin{multicols}{5}
            \begin{reponses}
                \mauvaise{8}
                \bonne{4}
                \mauvaise{5}
                \mauvaise{3}
                \mauvaise{Boucle infinie}
            \end{reponses}
        \end{multicols}
    \end{question}
}
\element{bank-boucle-offbyone}{
    \begin{question}{boucle-offbyone-v3}
        Combien de fois le message \texttt{"Bonjour"} est-il affiché par \lstinline[language=C]|for (int i = 1; i <= 12 && i != 4; i++) { printf("Bonjour"); }| ?
        \begin{multicols}{5}
            \begin{reponses}
                \mauvaise{12}
                \bonne{3}
                \mauvaise{4}
                \mauvaise{2}
                \mauvaise{Boucle infinie}
            \end{reponses}
        \end{multicols}
    \end{question}
}
\element{boucles}{
    \restituegroupe[1]{bank-boucle-offbyone}
}

% ---- boucle-ecrire (écriture de code : reste \evaluationProfGrille, pas un tableau) ----
\element{bank-boucle-ecrire}{
    \begin{question}{boucle-ecrire-v1}
        Écrire une boucle en langage C qui affiche les nombres entiers de 1 à 10 (une valeur par ligne).
        \evaluationProfGrille{2}{9}
    \end{question}
}
\element{bank-boucle-ecrire}{
    \begin{question}{boucle-ecrire-v2}
        Écrire une boucle en langage C qui affiche les nombres entiers de 1 à 6 (une valeur par ligne).
        \evaluationProfGrille{2}{9}
    \end{question}
}
\element{bank-boucle-ecrire}{
    \begin{question}{boucle-ecrire-v3}
        Écrire une boucle en langage C qui affiche les nombres entiers de 0 à 8 (une valeur par ligne).
        \evaluationProfGrille{2}{9}
    \end{question}
}
\element{boucles}{
    \restituegroupe[1]{bank-boucle-ecrire}
}
